\input{formattingHeader}

\titleandheader{Modular Transport Abstraction for the OpenLCB Checker}

\begin{document}
\maketitle
\thispagestyle{firststyle}

\section{Introduction}

This note documents a proposed design for making the OpenLCB Checker
transport-layer--agnostic, so that the same test suite can run over
CAN (GridConnect serial or TCP) or native OpenLCB TCP/IP, with a clean
extension point for future physical layers.

The design uses an abstract \texttt{TransportStack} base class with one
concrete implementation per physical layer.  A factory in \texttt{setup.py}
selects the correct stack based on a new \texttt{--transport} configuration
option.

\section{Current State}

All 93 \texttt{check\_*.py} tests currently run through a single code path in
\texttt{olcbchecker/setup.py} that hard-codes the CAN GridConnect stack:

\begin{verbatim}
CanPhysicalLayerGridConnect  ->  CanLink  ->  message listeners
\end{verbatim}

The underlying connection is either a TCP socket (GridConnect ASCII over TCP,
port 12021) or a serial port, both carrying CAN GridConnect frames.

A native TCP/IP link layer already exists in
\texttt{openlcb/tcplink/tcplink.py} and \texttt{openlcb/tcplink/tcpsocket.py}
but is not wired into the test infrastructure.

\subsection{Test Transport Dependency Analysis}

Of the 93 test files:

\begin{itemize}
\item \textbf{83 tests} use only \texttt{sendMessage()} and
    \texttt{getMessage()} --- already transport-agnostic.
\item \textbf{10 tests} use CAN-specific APIs and are inherently
    CAN-only:
\end{itemize}

\begin{tabular}{|l|p{8cm}|}
\hline
\textbf{Test File} & \textbf{CAN APIs Used} \\
\hline
\texttt{check\_fr10\_init.py} & \texttt{getFrame}, header bit operations \\
\hline
\texttt{check\_fr20\_ame.py} & \texttt{CanFrame}, \texttt{sendCanFrame}, \texttt{getFrame}, \texttt{canLink.localAlias} \\
\hline
\texttt{check\_fr30\_collide.py} & \texttt{CanFrame}, \texttt{sendCanFrame}, \texttt{getFrame}, \texttt{canLink.localAlias} \\
\hline
\texttt{check\_fr40\_highbit.py} & \texttt{CanFrame}, \texttt{sendCanFrame}, \texttt{getFrame}, \texttt{canLink.localAlias} \\
\hline
\texttt{check\_fr50\_capacity.py} & \texttt{canLink.localAlias} \\
\hline
\texttt{check\_fr60\_standard.py} & \texttt{CanFrame}, \texttt{getFrame}, \texttt{canLink.localAlias} \\
\hline
\texttt{check\_ev60\_ewp.py} & \texttt{CanFrame}, \texttt{sendCanFrame}, \texttt{canLink.localAlias} \\
\hline
\texttt{check\_ev70\_lrn.py} & \texttt{CanFrame}, \texttt{canLink.localAlias} \\
\hline
\texttt{check\_mc50\_restart.py} & \texttt{CanFrame}, \texttt{sendCanFrame}, \texttt{canLink.localAlias} \\
\hline
\texttt{check\_st110\_concurrent\_multi\_node.py} & \texttt{CanFrame}, \texttt{sendCanFrame}, \texttt{getFrame}, \texttt{canLink.nodeIdToAlias} \\
\hline
\end{tabular}

\section{Proposed Architecture}

\subsection{TransportStack Abstract Base Class}

A new file \texttt{olcbchecker/transportstack.py} defines the interface that
every transport must implement:

\begin{description}
\item[\texttt{link\_layer}] Returns the \texttt{LinkLayer} subclass instance
    (\texttt{CanLink}, \texttt{TcpLink}, etc.).
\item[\texttt{message\_queue}] A \texttt{Queue} of received
    \texttt{Message} objects.
\item[\texttt{send\_message(msg)}] Send an OpenLCB \texttt{Message}.
\item[\texttt{supports\_frames()}] Returns \texttt{True} if the transport
    exposes raw frame access (e.g.\ CAN), \texttt{False} otherwise.
\item[\texttt{frame\_queue}] A \texttt{Queue} of received frames, or
    \texttt{None} if the transport has no frame concept.
\item[\texttt{send\_frame(frame)}] Send a raw frame.  Raises
    \texttt{NotImplementedError} for transports without frames.
\item[\texttt{local\_alias}] The CAN alias of the local node, or
    \texttt{None} if not applicable.
\item[\texttt{node\_id\_to\_alias(node\_id)}] Resolve a NodeID to a CAN
    alias, or \texttt{None} if not applicable.
\item[\texttt{start()}] Start the background receive loop thread.
\item[\texttt{close()}] Close the underlying connection.
\end{description}

\subsection{CAN GridConnect Stack}

A new file \texttt{olcbchecker/can\_stack.py} containing class
\texttt{CanStack(TransportStack)}.

This extracts the existing logic from \texttt{setup.py}:
\begin{enumerate}
\item Creates \texttt{TcpSocket} (from \texttt{openlcb/canbus/tcpsocket.py})
    or \texttt{SerialLink} based on \texttt{hostname} vs.\ \texttt{devicename}.
\item Creates the \texttt{ReportingCanLink} subclass of
    \texttt{CanPhysicalLayerGridConnect}.
\item Creates \texttt{CanLink} with \texttt{localNodeID}.
\item Wires listeners: frames go to \texttt{frame\_queue}, messages go to
    \texttt{message\_queue}.
\item \texttt{supports\_frames()} returns \texttt{True}.
\item \texttt{local\_alias} and \texttt{node\_id\_to\_alias()} delegate to the
    underlying \texttt{CanLink}.
\item \texttt{start()} spawns the receive loop thread that calls
    \texttt{receiveString()}.
\end{enumerate}

\subsection{Native TCP Stack}

A new file \texttt{olcbchecker/tcp\_stack.py} containing class
\texttt{TcpStack(TransportStack)}.

\begin{enumerate}
\item Creates \texttt{TcpSocket} (from \texttt{openlcb/tcplink/tcpsocket.py})
    and connects to \texttt{hostname:portnumber}.
\item Creates \texttt{TcpLink} with \texttt{localNodeID}.
\item Wires \texttt{TcpLink.linkPhysicalLayer(tcpSocket.send)}.
\item Registers a message listener that puts received messages into
    \texttt{message\_queue}.
\item \texttt{supports\_frames()} returns \texttt{False}.
\item \texttt{frame\_queue} returns \texttt{None}.
\item \texttt{start()} spawns a receive loop that calls
    \texttt{tcpSocket.receive()} and feeds bytes to
    \texttt{tcpLink.receiveListener()}, then calls \texttt{tcpLink.linkUp()}.
\item \texttt{close()} calls \texttt{tcpLink.linkDown()} then
    \texttt{tcpSocket.close()}.
\end{enumerate}

\subsection{Factory in setup.py}

\texttt{olcbchecker/setup.py} is reduced to a thin factory:

\begin{verbatim}
if configure.transport == "tcp" :
    from olcbchecker.tcp_stack import TcpStack
    stack = TcpStack(configure)
else :
    from olcbchecker.can_stack import CanStack
    stack = CanStack(configure)

stack.start()
\end{verbatim}

Backward-compatible aliases (\texttt{canLink}, \texttt{sendCanFrame},
\texttt{frameQueue}, \texttt{messageQueue}) are set from the stack's
properties so that the 10 CAN-specific tests continue to work without
modification when running in CAN mode.

\subsection{Updates to olcbchecker/\_\_init\_\_.py}

\begin{itemize}
\item \texttt{sendMessage()} calls \texttt{setup.stack.send\_message()} instead
    of \texttt{setup.canLink.sendMessage()}.
\item New function \texttt{supports\_frames()} returns
    \texttt{setup.stack.supports\_frames()}.
\item \texttt{getFrame()} raises a clear \texttt{RuntimeError} when the
    transport does not support frames.
\item \texttt{purgeMessages()} only drains the frame queue when it exists.
\end{itemize}

\section{Configuration}

\subsection{New Configuration Variable}

\texttt{defaults.py} gains:

\begin{verbatim}
transport = "can"
# "can" for CAN GridConnect (TCP or serial)
# "tcp" for native OpenLCB TCP/IP
\end{verbatim}

This follows the existing pattern: \texttt{defaults.py} $\to$
\texttt{localoverrides.py} $\to$ command-line arguments.

\subsection{New Command-Line Option}

\texttt{configure.py} gains:

\begin{verbatim}
--transport TYPE    "can" (default) or "tcp"
\end{verbatim}

When \texttt{transport} is \texttt{"tcp"}, the checker requires
\texttt{hostname} to be set (native TCP has no serial mode).

\section{Test Gating for CAN-Only Tests}

\subsection{Frame Transport Suite}

\texttt{control\_frame.py} --- the entire Frame Transport suite (6 tests) is
CAN-only.  \texttt{checkAll()} checks \texttt{supports\_frames()} and returns
0 with an informational log message if false.  The menu in
\texttt{control\_master.py} hides option ``0 Frame Transport checking'' when
frames are not supported.

\subsection{Individual CAN-Specific Tests Outside Frame Suite}

Three CAN-specific tests live in other suites:

\begin{itemize}
\item \texttt{check\_ev60\_ewp.py} (Event Transport)
\item \texttt{check\_ev70\_lrn.py} (Event Transport)
\item \texttt{check\_mc50\_restart.py} (Memory Configuration)
\end{itemize}

And one in Stream Transport:

\begin{itemize}
\item \texttt{check\_st110\_concurrent\_multi\_node.py}
\end{itemize}

\textbf{Open question:} Should these be gated in their respective controllers
(\texttt{control\_events.py}, \texttt{control\_memory.py},
\texttt{control\_stream.py}), or should each test file check
\texttt{supports\_frames()} at the top of its \texttt{check()} function?
Gating in the controller is cleaner (no test file changes) but less
self-documenting.

\section{Adding a Future Physical Layer}

To add a new transport (e.g.\ USB, wireless):

\begin{enumerate}
\item Create one new file: \texttt{olcbchecker/usb\_stack.py} subclassing
    \texttt{TransportStack}.
\item Add one \texttt{elif} in \texttt{setup.py}'s factory.
\item Add the transport name to the \texttt{--transport} help text in
    \texttt{configure.py}.
\item If the transport has its own unique test category, add a
    \texttt{control\_usb.py} and wire it into \texttt{control\_master.py}.
\end{enumerate}

No other existing files require modification.

\section{TCP-Specific Tests (Future Scope)}

The native TCP specification defines link control messages (Status
Request/Reply, Drop Link Request/Reply) that could have their own test suite
analogous to the Frame Transport suite for CAN.  This would be a new
\texttt{control\_tcplink.py} with \texttt{check\_tl*.py} files, shown in the
menu only when \texttt{transport == "tcp"}.  This is out of scope for the
initial implementation but the architecture supports it.

\section{Files Changed Summary}

\subsection{New Files (3)}

\begin{tabular}{|l|p{9cm}|}
\hline
\textbf{File} & \textbf{Purpose} \\
\hline
\texttt{olcbchecker/transportstack.py} & Abstract base class defining the transport interface \\
\hline
\texttt{olcbchecker/can\_stack.py} & CAN GridConnect implementation (extracted from current \texttt{setup.py}) \\
\hline
\texttt{olcbchecker/tcp\_stack.py} & Native TCP/IP implementation (wires existing \texttt{TcpLink}) \\
\hline
\end{tabular}

\subsection{Modified Files (5)}

\begin{tabular}{|l|p{9cm}|}
\hline
\textbf{File} & \textbf{Change} \\
\hline
\texttt{defaults.py} & Add \texttt{transport = "can"} \\
\hline
\texttt{configure.py} & Add \texttt{--transport} CLI option and \texttt{localoverrides} chain entry \\
\hline
\texttt{olcbchecker/setup.py} & Replace hard-coded CAN stack with factory; add backward-compat aliases \\
\hline
\texttt{olcbchecker/\_\_init\_\_.py} & Use \texttt{stack.send\_message()}, add \texttt{supports\_frames()}, guard \texttt{getFrame()} \\
\hline
\texttt{control\_master.py} & Conditionally show/run Frame Transport suite based on \texttt{supports\_frames()} \\
\hline
\end{tabular}

\subsection{Unchanged (93 test files)}

No \texttt{check\_*.py} files are modified.  CAN-specific tests work
unchanged in CAN mode via backward-compatible aliases.  In TCP mode they are
skipped by the controllers.

\section{Open Questions}

\begin{enumerate}
\item \textbf{Port number default:} Should \texttt{transport = "tcp"}
    automatically switch the default port from 12021 to the native TCP port,
    or should the user always specify it explicitly?
\item \textbf{Gating CAN-specific tests outside the Frame suite:}
    Gate in the controller files, or put a \texttt{supports\_frames()} guard
    inside each of the 4 affected test files?
\item \textbf{CLI flag:} Is \texttt{--transport can|tcp} acceptable, or
    is a different naming preferred?
\end{enumerate}

\section{Verification}

\begin{enumerate}
\item Run with \texttt{--transport can -a localhost:12021} --- all 93 tests
    must behave identically to today.
\item Run with \texttt{--transport tcp -a host:port} against a native TCP
    device --- 83 message-level tests run; 10 CAN-specific tests are skipped
    with informational messages.
\item Menu shows ``Frame Transport checking'' only in CAN mode.
\item \texttt{checkAll()} reports correct section counts (skipped sections
    do not count as failures).
\item Running a CAN-specific test file directly on TCP gives a clear
    error, not a silent pass.
\end{enumerate}

\end{document}
